% GDZ PPN236018647,native590middle–591upper;HKUST591–592.
\NumberedItem{119}
\noindent\hspace*{\parindent}Methodus nova simplicissima expeditissima elementa orbitarum corporum coelestium investigandi.
\par\noindent\hfill Brunsv[igae, 1801] Sept. m[edio]
\par\medskip\begingroup\small
\renewcommand{\thefootnote}{*)}
\noindent\hspace*{\parindent}\textsc{Gauss} hatte sich schon seit längerer Zeit mit dem Studium der theoretischen Astronomie und nach Ausweis der Scheda Ac (begonnen Nov. 1799), S. 23, auch mit der \textsc{Olbers}schen Methode der Bestimmung parabolischer Bahnen\footnote{\textsc{W. Olbers}, \textit{Abhandlung über die leichteste und bequemste Methode, die Bahn eines Cometen zu berechnen}, Weimar 1797, Sein Leben und seine Werke I, Berlin 1894, S. 3.} beschäftigt, als der am 1. Januar 1801 von \textsc{Piazzi} in Palermo entdeckte erste kleine Planet \textls[120]{Ceres} bei seinem bevorstehenden Wiedererscheinen gegen Ende des Sommers 1801 von den Astronomen vergeblich am Himmel gesucht wurde. Die Schwierigkeit der Bahnbestimmung lag hauptsächlich darin, daß die Beobachtungen von \textsc{Piazzi} sich nur über den kurzen Zeitraum von 42 Tagen erstreckten, für welchen Fall die bis dahin üblichen Methoden nicht ausreichten. Die Aufgabe, durch genauere Bestimmung seiner Bahn die Wiederauffindung doch noch zu ermöglichen, übte auf \textsc{Gauss} eine starke Anziehungskraft aus und er verwandte von nun ab einen großen Teil seiner Arbeitszeit zu Untersuchungen und Rechnungen zur Bahnbestimmung, aus denen schließlich als vollendetes Werk die \textit{Theoria motus} (erschienen 1809) hervorging.
\par\noindent\hspace*{\parindent}Auf welche besondere Methode die vorstehende Eintragung Bezug nimmt, ob auf die in der Monatlichen Correspondenz zur Beförderung der Erd- und Himmelskunde, Juni 1802, abgedruckte\renewcommand{\thefootnote}{**)}\footnote{\textit{Vorschriften, um aus der geocentrischen Länge und Breite eines Himmelskörpers, dem Orte seines Knotens, der Neigung der Bahn, der Länge der Sonne und ihrem Abstände von der Erde abzuleiten: Des Himmelskörpers heliocentrische Länge in der Bahn, wahren Abstand von der Sonne und wahren Abstand von der Erde}, Werke VI, S. 87.} oder auf eine Vorgängerin dieser, läßt sich nicht entscheiden, da die ältesten im Nachlaß, namentlich in den Schedae Ag und Ah vorhandenen Aufzeichnungen und Rechnungen über diesen Gegenstand erst aus dem November 1801 stammen. Das gleiche gilt von der Tagebucheintragung Nr. 121.
\par\noindent\hspace*{\parindent}In der Vorrede zur \textit{Theoria motus} (Werke VII, 1906, S. 7) sagt \textsc{Gauss}:
\par\noindent\hspace*{\parindent}»satis mirum videtur, problema generale
\begin{quote}
Determinare orbitam corporis coelestis, absque omni suppositione hypothetica, ex observationibus tempus haud magnum complectentibus neque adeo delectum, pro applicatione methodorum specialium, patientibus
\end{quote}
\par\endgroup
\par\medskip\noindent\small X1.\hfill 71\normalsize

\clearpage\SourceLeaf{591}\ArithmeticHead{562}{L}{TAGEBUCH MIT ERLÄUTERUNGEN.}
\par\begingroup\small
\noindent usque ad initium huius saeculi penitus propemodum neglectum esse, vel saltem a nemine serio ac digne tractatum, quum certe theoreticis propter difficultatem atque elegantiam sese commendare potuisset, etiamsi apud practicos de summa eius utilitate nondum constaret. Scilicet invaluerat apud omnes opinio, \textls[120]{impossibilem} esse talem determinationem completam ex observationibus breviori temporis intervallo inclusis, male sane fundata, quum nunc quidem certissimo iam evictum sit, orbitam corporis coelestis ex observationibus bonis paucos tantummodo dies complectentibus, absque ulla suppositione hypothetica, satis approximate iam determinari posse.
\renewcommand{\thefootnote}{[*)]}
\par\noindent\hspace*{\parindent}Incideram in quasdam ideas, quae ad solutionem problematis magni de quo dixi facere videbantur, mense Septembri a. 1801, tunc in labore plane diverso occupatus\footnotemark.\quad .\ .\ .\ .\ .\ .\ .\ .\ .\ .\ .\ . eodem circiter tempore rumor de planeta novo Ian. 1. istius anni in specula Panormitana detecto per omnium ora volitabat, moxque ipsae observationes inde ab epocha illa usque ad 11. Febr. ab astronomo praestantissimo \textsc{Piazzi} institutae ad notitiam publicam pervenerunt.\quad .\ .\ .\ .\ .\ .\ .\ .\ .\ .\ .\ . Unquamne opportunius experiri potuissem, ecquid valeant ideolae meae ad usum practicum, quam si tunc istis ad determinationem orbitae Cereris uterer, qui planeta inter 41 illos dies geocentrice arcum trium tantummodo graduum descripserat, et post annum elapsum in coeli plaga longissime illinc remota indagari debebat? Prima haecce methodi applicatio facta est mense Oct. 1801, primaque nox serena, ubi planeta ad normam numerorum inde deductorum quaesitus est (Dec. 7, 1801 a clar. \textsc{de Zach}) transfugam observationibus reddidit«.
\renewcommand{\thefootnote}{[*)}
\footnotetext{Vergl. die folgende Nr. 120.]}
\par\noindent\hspace*{\parindent}Siehe auch die \textit{Anzeige} der \textit{Theoria motus}, Werke VI, S. 56 und den Brief an \textsc{Olbers} vom 25. Mai 1802, \textit{Wilhelm Olbers, Sein Leben und seine Werke} II1, Berlin 1900, S. 48.
\par\noindent\hspace*{\parindent}Vergl. die Nummern 122, 125, 126, 127, 129.
\par\noindent\hfill\textsc{Brendel.}
\par\endgroup
